\section{A quasi-filtered category that is not total}\label{sec:twographs}

\subsection{Ordered two-graphs}

Let $S$ be a finite set. A \emph{graph} on $S$ is a function $b$ from the
two-element subsets of $S$ to $\F_2$. A \emph{two-graph} on $S$ is a function
$c$ from the three-element subsets of $S$ to $\F_2$ such that, for every
four-element subset $Q$ of $S$, the sum of $c$ over the four three-element
subsets of $Q$ is zero. (The triples with $c=1$ then form a two-graph in the
usual sense; see \cite{Seidel}.) We write $c(x,y,z)$ for $c(\{x,y,z\})$ and
put $c(x,y,z)=0$ when two of the entries coincide; similarly for graphs. For
a graph $b$ let
\[
 (\partial b)(x,y,z)=b(x,y)+b(x,z)+b(y,z).
\]
This is a two-graph, because in the sum over the four triples of a
four-element set every pair occurs twice. For a two-graph $c$ on $S$ and
$p\in S$ let $c_p$ be the graph $c_p(x,y)=c(p,x,y)$. Then
\begin{equation}\label{eq:anchor}
 \partial(c_p)=c ;
\end{equation}
for a triple that does not contain $p$ this is the condition on the
four-element set formed with $p$, and for a triple $\{p,y,z\}$ it holds because
$c_p(p,y)=c_p(p,z)=0$.

An \emph{ordered two-graph} is a finite non-empty linearly ordered set with a
two-graph. An \emph{embedding} of ordered two-graphs is an increasing
injection $f\colon X\to Y$ with $c_Y(f(x),f(y),f(z))=c_X(x,y,z)$ for all
$x,y,z$.

\begin{definition}
The objects of $\mathcal T$ are the ordered two-graphs with underlying ordered
set $[n]$, $n\ge1$; its arrows are the embeddings.
\end{definition}

\begin{lemma}[anchored amalgams]\label{lem:anchored}
Let $S=S_1\cup S_2$ and let $c_i$ be a two-graph on $S_i$, $i=1,2$, such that
$c_1$ and $c_2$ agree on the triples contained in $S_1\cap S_2$. Choose
$p_i\in S_i$, with $p_1=p_2\in S_1\cap S_2$ if this set is non-empty, and let
$b$ be the graph on $S$ that equals $(c_i)_{p_i}$ on the pairs contained in
$S_i$ and is zero on all other pairs. Then $\partial b$ is a two-graph on $S$
that restricts to $c_i$ on $S_i$.
\end{lemma}

\begin{proof}
The graph $b$ is well defined, because $(c_1)_p$ and $(c_2)_p$ agree on the
pairs contained in $S_1\cap S_2$ when $p$ lies in this set. On triples
contained in $S_i$ we have $\partial b=\partial((c_i)_{p_i})=c_i$ by
\eqref{eq:anchor}.
\end{proof}

We call $\partial b$ the amalgam \emph{anchored at} $p_1,p_2$. In particular
any two objects of $\mathcal T$ embed into a third, so $\mathcal T$ is
connected.

\subsection{The generic ordered two-graph}

\begin{proposition}\label{prop:generic2}
There is an ordered two-graph structure on a countably infinite set $M$
\textup{(}a linear order and a function $c$ on three-element subsets that is a
two-graph on every finite subset\textup{)} such that:
\begin{enumerate}
\item[\textup{(N2)}] if $Y\subseteq M$ is finite and $f\colon Y\to C$ is an
 embedding into an ordered two-graph, there is an embedding $h\colon C\to M$
 with $hf=\id_Y$ \textup{(}here $Y$ may be empty\textup{)};
\item[\textup{(N3)}] every isomorphism between finite subsets of $M$ extends to an
 automorphism of $M$.
\end{enumerate}
\end{proposition}

\begin{proof}
As for Proposition~\ref{prop:generic}, with every finite subset in the role
of a closed subset and with the anchored amalgams of Lemma~\ref{lem:anchored},
ordered by any linear order extending the two given ones, in the role of the
free amalgams.
\end{proof}

Fix such an $M$, let $G=\Aut(M)$ and let $\K$ be the set of all finite subsets
of $M$. By (N3) the arrows of $\mathcal A=\mathcal A(M,G,\K)$ are the embeddings
between finite non-empty subsets of $M$, and by (N2) with $Y=\varnothing$ the
category $\mathcal A$ is equivalent to $\mathcal T$. Condition (K1) is clear.
By (N2) the order of $M$ has no greatest element.

\begin{lemma}\label{lem:K24two}
$(M,G,\K)$ satisfies \textup{(K2)} and \textup{(K4)}.
\end{lemma}

\begin{proof}
(K2). Let $B\subseteq M$ be finite and $x\in M\setminus B$. Let
$C=B\cup\{x,x'\}$ carry an amalgam of $B\cup\{x\}$ and of a second copy
$B\cup\{x'\}$ of it over $B$ (Lemma~\ref{lem:anchored}), ordered with $x'$
immediately after $x$. By (N2) it embeds into $M$ over $B\cup\{x\}$, and by
(N3) the isomorphism $B\cup\{x\}\to B\cup\{x'\}$ extends to some $g\in G_B$
with $gx\ne x$.

(K4). We argue as for Lemma~\ref{lem:SCdelta}. Let $\mathcal C$ be the
category of ordered two-graphs whose underlying set is a finite subset of
$\Omega$, with the embeddings as arrows; it is equivalent to $\mathcal A$. Let
$P$ be an object, $p_0$ its least element, $\mathcal U$ a non-empty cosieve of
$P\backslash\mathcal C$ and $(Z,z)$ an object of $\mathcal U$. For an object
$(X,x)$ let $T(X)$ be the set $\lambda(X)\cup\rho(Z\setminus z(P))$, ordered
as in Lemma~\ref{lem:SCdelta}, with the amalgam of the copies of $X$ and of
$Z$ anchored at $\lambda(x(p_0))$. For an arrow $f\colon X\to X'$ under $P$,
the map $T(f)$ defined as before is increasing. It carries the anchor to the
anchor and the pairs contained in one of the two copies to such pairs,
preserving the values of the graph $b$ of Lemma~\ref{lem:anchored}, and it
carries the remaining pairs, on which $b=0$, to such pairs. Hence it
preserves $\partial b$ and is an embedding. As before we obtain an
endofunctor $T$ of $P\backslash\mathcal C$ and natural transformations
$\id\to T\leftarrow\mathrm{const}_{(Z,z)}$, and $B\mathcal U$ is contractible.
\end{proof}

\subsection{Open subgroups}

\begin{lemma}\label{lem:stab2a}
Let $B,C\subseteq M$ be finite with $B\cap C\ne\varnothing$. Then
$\langle G_B,G_C\rangle=G_{B\cap C}$.
\end{lemma}

\begin{proof}
Let $g\in G_{B\cap C}$. Put $D=B\cup C$, $E=D\cup gD$, $F=B\cap C$, and let
$d_0\colon D\to E$ be the inclusion and $d_1=g|_D$. Let $Z$, $e$,
$a_0,\dots,a_s$ and the pieces $Z_i$ be as in Lemma~\ref{lem:bridgeset}. Each
piece is a copy of $E$ or of $D$ and carries the corresponding two-graph.
Two consecutive pieces induce the same two-graph on their intersection, which
is the common image of $B$ or of $C$, because $d_0$ and $d_1$ are embeddings;
call this $(\ast)$. Fix $p\in F$ and let $\hat p$ be the corresponding
$F$-point of $Z$. Let $b$ be the graph on $Z$ with
$b(u,u')=c_{Z_i}(\hat p,u,u')$ if $u$ and $u'$ lie in a common piece $Z_i$, and
$b(u,u')=0$ if they lie in no common piece. This is well defined. Indeed, let
$u,u'\in Z_i\cap Z_j$. If $i$ and $j$ are consecutive, apply $(\ast)$. If not,
and $u$, $u'$ are both $F$-points, note that all pieces induce the same
two-graph on the set of $F$-points. Otherwise $\{i,j\}=\{1,s-1\}$ and
$u,u'\in Z_0$ by Lemma~\ref{lem:bridgeset}(ii), and we apply $(\ast)$ to
$(Z_j,Z_0)$ and to $(Z_0,Z_i)$. Give $Z$ the two-graph $\partial b$. On a
triple contained in a piece $Z_i$ it equals $\partial((c_{Z_i})_{\hat p})=c_{Z_i}$
by \eqref{eq:anchor}. So $e$ and all $a_i$ are embeddings of ordered
two-graphs.

By (N2) there is an embedding $h\colon Z\to M$ with $he=\id_E$, and by (N3)
$ha_i=g_i|_D$ with $g_i\in G$, $g_0=1$, $g_s=g$. By
Lemma~\ref{lem:bridgeset}(i), $g_i^{-1}g_{i+1}$ lies in $G_B$ for $i$ even and
in $G_C$ for $i$ odd, so $g\in\langle G_B,G_C\rangle$.
\end{proof}

\begin{lemma}\label{lem:stab2b}
For distinct points $x,y$ of $M$ we have $\langle G_{\{x\}},G_{\{y\}}\rangle=G$.
\end{lemma}

\begin{proof}
Let $H$ be the subgroup on the left, and suppose $x<y$. A two-graph on a set
with at most two elements is empty, so by (N3) the group $G$ is transitive on
$M$, and $G_{\{w\}}$ is transitive on $\{u:u<w\}$ and on $\{u:u>w\}$ for every
$w\in M$. Note that $H$ contains $G_{\{hx\}}$ and $G_{\{hy\}}$ for every
$h\in H$. Let $u\in M$ and choose $w$ greater than $u$ and $y$. Then
$w\in G_{\{x\}}\,y$, so $G_{\{w\}}\le H$, and $u\in G_{\{w\}}\,x$. Hence $H$
is transitive on $M$. For $g\in G$ choose $h\in H$ with $hx=gx$; then
$h^{-1}g\in G_{\{x\}}\le H$, so $g\in H$.
\end{proof}

\begin{proposition}\label{prop:open2}
For all finite $B,C\subseteq M$ we have $\langle G_B,G_C\rangle=G_{B\cap C}$, and
$(M,G,\K)$ satisfies \textup{(K3)}.
\end{proposition}

\begin{proof}
If $B\cap C\ne\varnothing$ this is Lemma~\ref{lem:stab2a}. Otherwise choose
distinct points $x,y$ outside $B\cup C$. By Lemma~\ref{lem:stab2a},
$\langle G_B,G_C\rangle$ contains
$\langle G_{B\cup\{x\}},G_{C\cup\{x\}}\rangle=G_{\{x\}}$, and likewise
$G_{\{y\}}$; so it is $G$ by Lemma~\ref{lem:stab2b}. Condition (K3) follows as
in Corollary~\ref{cor:opendelta}.
\end{proof}

\subsection{The classifying space}

\begin{lemma}\label{lem:pi1}
$B\mathcal T$ is simply connected.
\end{lemma}

\begin{proof}
$\mathcal T$ is connected, so by the correspondence between covering spaces
and functors that invert all arrows (Section~\ref{sec:classes}) it
suffices to show that every functor $\Psi\colon\mathcal T\to\Set$ that sends
all arrows to bijections is isomorphic to a constant functor. Let $P$ and
$Q$ be the objects with one and with two points (there are no others with at
most two points), and $d^1,d^0\colon P\to Q$ the arrows onto the first and the
second point. Let $X$ be an object with three points and
$e_{ij}\colon Q\to X$ the embedding onto $\{i,j\}$. Then
\[
 e_{12}d^1=e_{13}d^1,\qquad e_{13}d^0=e_{23}d^0,\qquad e_{12}d^0=e_{23}d^1 .
\]
Since $\Psi$ takes values in bijections, the first two equations give
$\Psi(e_{12})=\Psi(e_{13})=\Psi(e_{23})$, and then the third gives
$\Psi(d^0)=\Psi(d^1)$. For an object $Y$ and $u\in Y$ let $i_u\colon P\to Y$ be
the arrow with image $\{u\}$. If $u<u'$ and $\epsilon\colon Q\to Y$ is the
embedding onto $\{u,u'\}$, then $i_u=\epsilon d^1$ and $i_{u'}=\epsilon d^0$, so
$\Psi(i_u)=\Psi(i_{u'})$. Thus $\theta_Y=\Psi(i_u)\colon\Psi(P)\to\Psi(Y)$ does
not depend on $u$. For an arrow $f\colon Y\to Y'$ we have $fi_u=i_{f(u)}$,
hence $\Psi(f)\theta_Y=\theta_{Y'}$. So $\theta$ is a natural isomorphism from
the constant functor with value $\Psi(P)$ to $\Psi$.
\end{proof}

\begin{lemma}\label{lem:H2}
$H_2(B\mathcal T;\F_2)\ne 0$.
\end{lemma}

\begin{proof}
We use the chains of the nerve with coefficients in $\F_2$. For a
$2$-simplex $\sigma=(X_0\xrightarrow{f_1}X_1\xrightarrow{f_2}X_2)$ let
$m_0\le m_1\le m_2$ be the images in $X_2$ of the greatest elements of $X_0$,
$X_1$, $X_2$, and put $\Phi(\sigma)=c_{X_2}(m_0,m_1,m_2)$. Then $\Phi$ is a
cocycle. Indeed, for a $3$-simplex $X_0\to X_1\to X_2\to X_3$ let
$n_0\le n_1\le n_2\le n_3$ be the images in $X_3$ of the greatest elements.
Since $X_2\to X_3$ is an embedding, the values of $\Phi$ on the four faces are
the values of $c_{X_3}$ on the four triples obtained by omitting one of the
$n_k$. If the $n_k$ are distinct, their sum is zero by the definition of a
two-graph. If $n_k=n_{k+1}$ for some $k$, the triples obtained by omitting
$n_k$ or $n_{k+1}$ give the same value, and the other two contain a
repetition.

Let $X$ and $X'$ be the two objects with three points, with
$c_X(1,2,3)=1$ and $c_{X'}(1,2,3)=0$. For $S_0\subset S_1\subset[3]$ with
$|S_0|=1$, $|S_1|=2$ let $\sigma(S_0,S_1,X)$ be the $2$-simplex
$P\to Q\to X$ in which the image of $Q$ is $S_1$ and the image of $P$ is
$S_0$; define $\sigma(S_0,S_1,X')$ likewise. Let $\xi$ be the sum of these
twelve simplices. It is a cycle: the face $P\to Q$ does not depend on the
target $X$ or $X'$; the face $Q\to X$ depends only on $S_1$ and so occurs
twice, and the face $P\to X$ depends only on $S_0$ and occurs twice; similarly
for $X'$. By definition $\Phi(\sigma(S_0,S_1,X))=c_X(\max S_0,\max S_1,3)$,
which is non-zero only for $S_0=\{1\}$, $S_1=\{1,2\}$, and $\Phi$ vanishes on
the simplices with target $X'$. So $\Phi(\xi)=1$, and $\xi$ is not a boundary.
\end{proof}

\begin{theorem}\label{thm:twographs}
Every action component of $\mathcal T$ is weakly contractible or weakly
equivalent to $\mathcal T$. Consequently $\mathcal T$ is quasi-filtered, and
it is not total.
\end{theorem}

\begin{proof}
By Lemma~\ref{lem:K24two} and Proposition~\ref{prop:open2} the triple
$(M,G,\K)$ satisfies (K1)--(K4), so the first assertion holds for $\mathcal A$
by Proposition~\ref{prop:reduction}, and hence for the equivalent category
$\mathcal T$. By Lemma~\ref{lem:pi1} all action components of $\mathcal T$ are
simply connected (each is weakly contractible or weakly equivalent to
$\mathcal T$), so $\mathcal T$ is quasi-filtered by
Proposition~\ref{prop:qf}. The category of elements of the functor with value
a point is $\mathcal T$ itself, which is not weakly contractible by
Lemma~\ref{lem:H2}; so $\mathcal T$ is not total by
Proposition~\ref{prop:total}.
\end{proof}

\begin{remark}\label{rem:tower}
The run from which this example is taken identifies $B\mathcal T$ with an
Eilenberg--Mac~Lane space $K(\F_2,2)$, and treats $n$-cocycles in the same way
for every $n\ge2$; see Section~\ref{sec:provenance}. Those statements are not
used and not claimed here.
\end{remark}
